\section{$c=0$ の層の完全な被約分類}
\label{sec:c-zero-v33}

この節では重みを $(1,1,1,3)$ に固定する。Poisson スキームを
$\mathcal P$ とし、$c=0$ で切り出した閉部分スキームを
$\mathcal Z$ と書く。以下は $\mathcal Z_{\mathrm{red}}$ の分類であり、
$\mathcal Z$ の冪零構造を除去してよいという主張ではない。

\begin{theorem}[零三次形式の層]
\label{thm:c-zero-v33}
$\mathcal Z_{\mathrm{red}}$ は、次元 $35,20,18$ のちょうど三つの
既約成分を持つ。これらはそれぞれ、次の三族の閉包である。
\begin{enumerate}
\item $D=0$、$[F]\in S_4^{\oplus3}/xS_3$ は任意。
\item $D\ne0$ の最小多項式の次数が $2$ で、$[F]$ は対応する
十五次元の混合核を動く。
\item $D$ の最小多項式の次数が $3$ で、$[F]=[gD]$、$g\in S_3$。
\end{enumerate}
これらの族は、全 Poisson スキームにおいてそれぞれ既知の成分
\[
 \operatorname{PB}(\mathbb P^2),\qquad
 L(1,1,4),\qquad L(1,1,1,3)
\]
に含まれる。従って $c=0$ に台を持つ未知の既約成分は存在しない。
\end{theorem}

\begin{proof}
$c=0$ では残る方程式は
\begin{equation}
 (x\times D)\cdot F=0.
 \label{eq:c-zero-det-v33}
\end{equation}
まず $D$ の最小多項式が三次とする。このとき各固有空間は一次元であり、
射影空間上の $x\times D=0$ の零点は有限個である。従って三つの二次式
$x\times D$ が生成するイデアルの高さは二である。
行列 $[\,x\mid D\,]$ の最大小行列式は、符号を除きこの三式である。
商体上で $x,D$ は一次独立なので、式 \eqref{eq:c-zero-det-v33} の解は
$F=fx+gD$ と一意に書ける。高さ一の素イデアルで局所化すれば、
上の最大小行列式の少なくとも一つが可逆であるため、二行を使って解くことで
$f,g$ はその局所環に属する。多項式環は正規環であり、高さ一の局所環の共通部分に
等しいので、$f,g$ は多項式である。斉次性から次数はいずれも三となる。
これは本例の Hilbert--Burch 完全性を直接示す議論である。従って全解は一意に
\[
 F=fx+gD,\qquad f,g\in S_3
\]
と書ける。Euler 商では $[F]=[gD]$ となり、核の次元は十である。
ここで $x,D$ は商体上で一次独立なので、$g\mapsto[gD]$ は単射である。

相異なる固有値を持つ $D=\operatorname{diag}(\lambda_1,\lambda_2,\lambda_3)$
の場合、$u=w-g$ と置くと、対応する水平一形式は
\[
 \alpha=(w-g)(x\times D)\cdot dx
 =xyzu\left((\lambda_3-\lambda_2)\frac{dx}{x}
 +(\lambda_1-\lambda_3)\frac{dy}{y}
 +(\lambda_2-\lambda_1)\frac{dz}{z}\right).
\]
これは $u$ に沿う留数を零とした $L(1,1,1,3)$ の一形式である。
三次の最小多項式を持つ $D$ の集合は $\mathfrak{sl}_3$ の既約な開集合であり、
その上の全解は上記の十次元ベクトル束である。
相異なる固有値の集合は稠密なので、その全体が同じ成分の閉包に含まれる。
この族の次元は $8+10=18$ である。

次に $D\ne0$ の最小多項式の次数が二の場合、Jordan 標準形により、
座標変更と非零の定数倍を除いて次の二型のみである。
\[
 D=\lambda(x\partial_x+y\partial_y-2z\partial_z)\quad(\lambda\ne0),
 \qquad D=x\partial_z.
\]
いずれも式 \eqref{eq:c-zero-det-v33} は $xF_2=yF_1$ と同値である。
従って $F_1=xh,F_2=yh$ と書け、Euler 商で一意に
\[
 F=(0,0,f),\qquad f\in S_4
\]
とできる。この核は十五次元である。
二型に対応する一形式はそれぞれ
\[
 (3\lambda zw+f)(x\,dy-y\,dx),\qquad
 (-xw+f)(x\,dy-y\,dx)
\]
である。いずれも $Q(x\,dy-y\,dx)$ という形を持ち、
因子 $x,y,Q$、重み $(1,1,4)$、留数 $(-1,1,0)$ の対数型に属する。
因子が特殊化している場合も $L(1,1,4)$ の閉包への包含が従う。

最小多項式の次数が二以下の行列の閉集合は既約で次元五である。
この集合は $I,D,D^2$ の一次従属という閉条件で定まる。
実際、階数一以下の行列 $N$ に対して
\[
 D=N-\frac{\operatorname{tr}N}{3}I
\]
とすると、$\operatorname{tr}N\ne0$ では二重固有値を持つ半単純型、
$\operatorname{tr}N=0$ では階数一の冪零型となり、上の非零二型を全て覆う。
階数一以下の行列の錐は既約で次元五であり、この写像は単射である。
単射性は、相異なる二つの階数一行列の差が非零スカラー行列にはなれないことから従う。
零行列を除いた基底上では混合核の次元は一定の十五であり、
対応する族はベクトル束、従って既約で次元 $5+15=20$ である。

最後に $D=0$ のファイバーは全 $35$ 次元の線形空間であり、
既知の $\operatorname{PB}(\mathbb P^2)$ 成分そのものである。
三つの局所閉族は $\mathcal Z$ の全ての閉点を尽くす。
十八次元族の一般点では $D$ は最小多項式三次なので、二十次元族の閉包に含まれない。
二十次元族は次元から十八次元族の閉包に含まれない。
また非零の $D$ を持つこれらの族は $D=0$ の三十五次元空間には含まれず、
逆向きの包含も次元から不可能である。
従って三つの閉包が、ちょうど $\mathcal Z_{\mathrm{red}}$ の既約成分である。

全 Poisson スキームの既約成分が $c=0$ に含まれていれば、
その既約性から上記三つの閉包のいずれかに含まれる。
いずれも既知の全スキームの既約成分に含まれるため、
成分の極大性から未知の成分は生じない。
\end{proof}

\begin{remark}[一般重みでの零三次形式の層]
同じ議論は重み $(1,1,1,m)$、$m\ge3$ にも適用できる。
この場合の混合項は $F\in S_{m+1}^{\oplus3}/xS_m$ であり、
$c=0$ の方程式は非零定数倍を除き依然として $\det(x,D,F)=0$ である。
最小多項式三次の型では $[F]=[gD]$、$g\in S_m$、
非零の二次型では $[F]=[(0,0,f)]$、$f\in S_{m+1}$ となる。
従って零三次形式の層の被約化は、ちょうど三つの既約成分を持ち、
それらの次元はそれぞれ
\[
 m^2+6m+8,\qquad 5+\binom{m+3}{2},\qquad
 8+\binom{m+2}{2}
\]
である。証明は上の次数四・三を次数 $m+1$・$m$ に置き換えるだけである。

この重みでの水平形式は、$c=0$ では
$\alpha=\tfrac m3w(x\times D)\cdot dx+(F\times x)\cdot dx$ となる。
従って後二族はそれぞれ対数型の族
$L(1,1,m+1)$ と $L(1,1,1,m)$ の閉包に含まれる。
ここでは、一般の $m$ においてこれらの対数型の閉包が全 Poisson スキームの
既約成分であることは主張しない。
定理~\ref{thm:c-zero-v33}の「$c=0$ に台を持つ未知の全スキームの成分はない」
という結論は、包含先の成分性が確立している $m=3$ に限って用いる。
\end{remark}

この定理は、$c\ne0$ の非被約三次形式や特殊 Jordan 型に関する残る分類を
自動的に解決するものではない。また $\mathcal Z_{\mathrm{red}}$ の
次元 $20,18$ の成分は、全 Poisson スキームの既約成分ではない。
